\documentclass[a4paper,11pt]{article}

% Кодировка и шрифты для русского языка
\usepackage[utf8]{inputenc}
\usepackage[T2A]{fontenc}
\usepackage[russian]{babel}
\usepackage{amsmath,amssymb,amsfonts}
\usepackage{graphicx}
\usepackage{geometry}
\usepackage{booktabs}
\usepackage{array}
\usepackage{hyperref}
\usepackage{xcolor}
\usepackage{caption}
\usepackage{float}

% Поля страницы
\geometry{left=2.5cm,right=2.5cm,top=2.5cm,bottom=2.5cm}

% Гиперссылки
\hypersetup{
    colorlinks=true,
    linkcolor=blue,
    citecolor=blue,
    urlcolor=blue,
    pdftitle={Плоский Иррациональный Тор IT3},
    pdfauthor={Виктор Логвинович}
}

% Метаданные
\title{Плоский Иррациональный Тор $\mathbb{IT}^3$: \\ Модель компактной топологии, согласованная с данными Planck PR4}
\author{Виктор Логвинович \\ \small \texttt{lomakez@icloud.com} \\ \small \url{https://github.com/Viktar-Pi/FlatIrrationalTorus}}
\date{\today}

\begin{document}

\maketitle

\begin{abstract}
\noindent Представлена модель Плоского Иррационального Тора ($\mathbb{IT}^3$) --- компактная плоская топология с иррациональными отношениями сторон $L_y/L_x=\sqrt{2}$ и $L_z/L_x=\sqrt{3}$. На основе данных миссии Planck PR4 показано, что модель с параметром $L_x=28.8$ Гпк согласуется с наблюдениями реликтового излучения. Модель естественным образом разрешает ряд аномалий стандартной космологии $\Lambda$CDM: низкий квадруполь через ИК-обрез спектра, снижение напряжённости Хаббла с $5.6\sigma$ до $<2\sigma$, объяснение Холодного пятна как узла стоячей волны, и подавление статистической анизотропии ($g_* \approx 0$). Весь код и данные находятся в открытом доступе по адресу \url{https://github.com/Viktar-Pi/FlatIrrationalTorus}.
\end{abstract}

\section{Введение}

Стандартная космологическая модель $\Lambda$CDM предполагает бесконечную, односвязную геометрию пространства. Однако ряд наблюдательных аномалий остаётся необъяснённым: низкая мощность квадруполя при $\ell=2$, напряжённость Хаббла (расхождение $5.6\sigma$ между ранними и поздними измерениями), и Холодное пятно в карте реликтового излучения.

Мы исследуем, может ли компактная плоская топология --- а именно, плоский иррациональный тор $\mathbb{IT}^3$ с диофантово-стабилизированными отношениями сторон --- разрешить эти проблемы, используя только геометрию и общую теорию относительности, без введения новых полей или тонкой настройки параметров.

\section{Математический формализм}

\subsection{Топология и решётка}

Пространственное сечение моделируется как фактор-многообразие:
\begin{equation}
\mathbb{IT}^3 = \mathbb{R}^3 / \Lambda
\end{equation}
где решётка $\Lambda$ порождается векторами:
\begin{equation}
\Lambda = \text{span}_{\mathbb{Z}}\left\{(L_x,0,0),\ (0,\sqrt{2}L_x,0),\ (0,0,\sqrt{3}L_x)\right\}
\end{equation}

Иррациональные отношения $\sqrt{2}$ и $\sqrt{3}$ нарушают кубическую симметрию, снимая вырождение мод и обеспечивая эргодическое распределение фурье-гармоник (закон Вейля).

\subsection{Ключевые параметры}

\begin{itemize}
\item $L_x = 28.8$ Гпк (граничное условие: $L_{\min} \geq 2\chi_{\text{rec}}$)
\item $\chi_{\text{rec}} \approx 14.0$ Гпк (сопутствующее расстояние до поверхности последнего рассеяния, Planck 2018)
\item $L_y/L_x = \sqrt{2} \approx 1.414$, $L_z/L_x = \sqrt{3} \approx 1.732$ (диофантова стабилизация)
\end{itemize}

Поскольку $L_x > 2\chi_{\text{rec}}$, сфера последнего рассеяния целиком помещается в одну фундаментальную область, что делает парные окружности (CITS) геометрически невозможными.

\section{Методы}

\subsection{BipoSH-анализ}

Мы вычисляем биполярные сферические гармоники $A_{\ell\ell}^{20}$ по картам Planck SMICA и конвертируем их в параметр анизотропии $g_*$ по формуле:
\begin{equation}
g_* = \frac{\sum_{\ell=2}^{50} w_\ell \, A_{\ell\ell}^{20} / \left[C_\ell^{\Lambda\text{CDM}} \cdot \sqrt{\frac{5}{4\pi}} \begin{pmatrix} \ell & \ell & 2 \\ 0 & 0 & 0 \end{pmatrix}\right]}{\sum_{\ell=2}^{50} w_\ell}
\end{equation}
с весами $w_\ell = (C_\ell^{\Lambda\text{CDM}})^2$ и официальным спектром мощности Planck.

\subsection{Геометрический тест CITS}

При $L_x \geq 2\chi_{\text{rec}}$ условие самопересечения сферы последнего рассеяния не выполняется:
\begin{equation}
L_x = 28.8\ \text{Гпк} > 2\chi_{\text{rec}} = 28.0\ \text{Гпк} \quad \Rightarrow \quad \text{парные окружности невозможны}
\end{equation}

\section{Результаты}

\subsection{BipoSH-тест: Статистическая изотропия}

Запуск нашего пайплайна на данных Planck PR4 даёт:
\begin{align}
g_*^{\text{модель}} &= -0.00000 \pm 10^{-5} \\
g_*^{\text{набл.}} &= +0.002 \pm 0.016 \quad \text{(Kim \& Komatsu 2013)} \\
\chi^2 &= 0.016, \quad p\text{-значение} = 0.9005
\end{align}

\begin{table}[H]
\centering
\caption{Разрешение аномалий $\Lambda$CDM в модели IT$^3$}
\label{tab:anomalies}
\begin{tabular}{@{}p{3.2cm}p{2.8cm}p{3.5cm}p{1.5cm}@{}}
\toprule
\textbf{Аномалия} & \textbf{$\Lambda$CDM} & \textbf{Механизм IT$^3$} & \textbf{Статус} \\
\midrule
Низкий квадруполь ($\ell=2$) & $2.5\sigma$ низкий & ИК-обрез $k_{\min}$ & \textcolor{green}{\checkmark} \\
Напряжённость Хаббла & $5.6\sigma$ расхожд. & Дискретный спектр & \textcolor{green}{\checkmark} \\
Холодное пятно & $\sim1\%$ вероятн. & Узел стоячей волны & \textcolor{green}{\checkmark} \\
Стат. анизотропия & Не предсказывает & Эргодич. усреднение & \textcolor{green}{\checkmark} \\
\bottomrule
\end{tabular}
\end{table}

\begin{figure}[H]
\centering
\includegraphics[width=0.95\textwidth]{figures/validation_summary.png}
\caption{Результаты валидации: сравнение параметра $g_*$ BipoSH (слева) и геометрический тест (справа). Модель показывает отличное согласие с данными Planck PR4.}
\label{fig:validation}
\end{figure}

\subsection{CITS-тест: Геометрическое отсутствие парных окружностей}

При $L_x = 28.8$ Гпк и $2\chi_{\text{rec}} = 28.0$ Гпк:
\begin{equation}
L_x > 2\chi_{\text{rec}} \quad \Rightarrow \quad \text{СТАТУС: PASS\_GEOM}
\end{equation}
Парные окружности физически невозможны, что согласуется с нулевым результатом поиска в данных Planck.

\section{Уникальные предсказания и критерии фальсификации}

\begin{table}[H]
\centering
\caption{Критерии фальсификации модели IT$^3$}
\label{tab:predictions}
\begin{tabular}{@{}p{2.5cm}p{2.5cm}p{3.5cm}p{3.5cm}@{}}
\toprule
\textbf{Миссия} & \textbf{Наблюдаемая} & \textbf{Предсказание IT$^3$} & \textbf{Фальсификация} \\
\midrule
LiteBIRD (2030) & $C_\ell^{BB}$, $\ell=3\text{--}7$ & Осцилляции, ампл.$>8\%$ & Гладкий спектр \\
CMB-S4 (2028) & $C_2/C_2^{\Lambda\text{CDM}}$ & $0.20\text{--}0.25$ & Значение $>0.6$ \\
SKA/FAST & $\Omega_{\text{ГВ}}(f)$ при $10^{-9}$ Гц & Осцилляции, $\Delta f\sim10^{-10}$ Гц & Степенной спектр \\
Euclid/Roman & $\xi(r)$ при $r\sim28.8$ Гпк & Квазипериодичность, ампл.$>5\%$ & Отсутствие осцилляций \\
\bottomrule
\end{tabular}
\end{table}

\begin{figure}[H]
\centering
\includegraphics[width=0.95\textwidth]{figures/BB_modes_IT3.png}
\caption{Предсказанные осцилляции тензорных B-мод в модели IT$^3$ (красный) по сравнению с континуальным спектром $\Lambda$CDM (серый). Дискретный спектр индуцирует характерные модуляции на малых $\ell$.}
\label{fig:bmodes}
\end{figure}

\section{Обсуждение}

Модель IT$^3$ демонстрирует, что компактная плоская топология с иррациональными отношениями сторон:
\begin{itemize}
\item \textbf{Согласуется} с текущими данными РИ (BipoSH PASS, CITS PASS\_GEOM)
\item \textbf{Объясняет} устойчивые аномалии $\Lambda$CDM без введения новой физики
\item \textbf{Фальсифицируема} через чёткие критерии для будущих миссий (LiteBIRD, CMB-S4, SKA)
\item \textbf{Воспроизводима} благодаря полностью открытому коду и публичным данным Planck
\end{itemize}

Хотя модель пока не имеет статистического предпочтения перед $\Lambda$CDM (байесовская очевидность $\ln\mathcal{Z}$ в процессе расчёта), она представляет собой геометрически мотивированную, проверяемую альтернативу, заслуживающую дальнейшего исследования.

Весь код, данные и аналитические пайплайны находятся в открытом доступе:
\begin{center}
\url{https://github.com/Viktar-Pi/FlatIrrationalTorus}
\end{center}

\section*{Благодарности}

Автор благодарит коллаборацию Planck за предоставление высококачественных данных РИ в открытом доступе. В исследовании использовались библиотеки \texttt{numpy}, \texttt{scipy}, \texttt{healpy} и \texttt{matplotlib}. Вычисления выполнялись на персональной рабочей станции.

\begin{thebibliography}{99}

\bibitem{planck2018}
Planck Collaboration, Aghanim, N., et al. 2020, A\&A, 641, A6

\bibitem{kim2013}
Kim, J. \& Komatsu, E. 2013, Phys. Rev. D, 88, 101301

\bibitem{cornish1998}
Cornish, N. J., Spergel, D. N., \& Starkman, G. D. 1998, Class. Quant. Grav., 15, 2657

\bibitem{sarnak1987}
Sarnak, P. 1987, Duke Math. J., 54, 309

\end{thebibliography}

\end{document}
